\section{Conclusion}
This paper frames wood species identification as a data governance problem rather than solely a model optimization problem. Instead of treating data partitioning as a post-collection detail, SCDP formalizes provenance tracking, hierarchical storage, manifest management, data auditing, and group-disjoint allocation as integral methodological components. CEGS-Split implements the allocation phase by assigning all images of a specimen to a single partition, maintaining per-species class balance, and using frozen embeddings to flag cases requiring provenance review.

Experiments on S3 demonstrate that the same KNN pipeline produces substantially different results under random image split versus CEGS-Split. This does not establish that every image-level benchmark contains leakage, but it shows that the split unit is an assumption with direct consequences for model evaluation conclusions. Recognition, metric learning, and self-supervised learning results should therefore always be reported together with the definition of the source unit, the split manifest, and the model selection rule.

Future work includes releasing and versioning the manifest, generating multiple group-disjoint splits or employing group cross-validation when specimen counts permit, and constructing an out-of-domain test set with independent imaging equipment and specimen sources. Integrating data auditing with wood anatomy experts is also necessary to distinguish genuinely biological signals from signals introduced by the image acquisition pipeline.
